\chapter{Network Segmentation}
\label{ch:3.3}
\chaptertopic{Topic 3.3}

\chapterquote{He who defends everything defends nothing.}{attributed to Frederick the Great}

\section{Motivation}

Chapter~\ref{ch:3.1} noted, almost in passing, that a flat network gives a compromised
workstation direct visibility to every server, printer, and camera on it. This chapter takes that
observation seriously and builds it into a full defensive strategy: dividing a network into zones
so that a breach in one zone cannot automatically reach another. This chapter covers four
techniques that do that division, each enforced at a different layer. Routed subnets and VLANs are
the coarse-grained tools most organizations start with. Private VLANs and Zero Trust
microsegmentation narrow that boundary down to a single host. Physical air gaps remove the network
connection entirely. The chapter moves from the coarsest of these tools to the most demanding, and
closes with the one that removes the network altogether.

\begin{apalignment}{Unit 3: Securing Networks}{Topic 3.3 -- Protecting Networks: Segmentation}
\begin{apobjectives}
\apobjective{3.3.A}{Mitigate Risk}{Identify the techniques by which a network is divided into
separately governed zones -- routed subnets, VLANs and their private-VLAN refinement, host-level
microsegmentation under a Zero Trust policy, and physical air gaps -- and know at which layer each
is enforced.}
\apobjective{3.3.B}{Mitigate Risk}{Explain why dividing a network raises its security: what a
boundary denies an intruder who is already inside, how it shrinks the reach of a single compromise,
and where segmentation stops helping.}
\end{apobjectives}
\end{apalignment}

\section{VLANs and Segmentation}

The coarsest form of segmentation predates VLANs entirely: dividing a network into separate IP
subnets, each behind its own router interface, so that traffic between subnets must pass through
-- and can be filtered by -- a router. VLANs formalize that same layer-3 idea inside a single
physical switch, without requiring a separate cable plant for every subnet.

\begin{definition}{VLAN (Virtual Local Area Network)}
A logical broadcast domain created in switch software rather than by physical cabling, so that
traffic broadcast on one VLAN never reaches devices on another unless a router or layer-3 switch
explicitly forwards it.
\end{definition}

A single physical switch can host dozens of VLANs simultaneously. Think of VLANs as walls inside
one building: the structure is a single physical switch, but the walls separate HR from
Engineering from the guest lobby. A fire in the lobby does not automatically spread to the server
room, because the wall is in the way.

An administrator assigns each switch port to a VLAN by numeric ID. A device plugged into a
VLAN~10 port sends traffic that stays on VLAN~10. Reaching a device on VLAN~20 instead requires
that traffic to leave through a router -- a process called \textbf{inter-VLAN routing} -- and only
if a firewall rule explicitly permits it.

\begin{example}[Segmentation limiting a compromise]
An attacker compromises a guest laptop on the visitor Wi-Fi VLAN. Without segmentation, that laptop
could scan and reach every device on the network. With segmentation in place, the damage stays
confined to the guest VLAN specifically. There is no physical layer-2 path to the corporate file
server on the internal employee VLAN, and the router forwards only what an access-control list
explicitly permits between the two. \textbf{Lateral movement} -- an attacker's technique for
hopping from one compromised host toward a more valuable target -- hits a wall at the first VLAN
boundary instead of continuing unopposed.
\end{example}

A common small-to-medium design uses four VLANs. The \textbf{management VLAN} carries switches,
routers, and infrastructure devices, never exposed to ordinary users. The \textbf{internal employee
VLAN} carries staff workstations and printers. The \textbf{server VLAN} carries databases and file
shares. The \textbf{guest/IoT VLAN} carries visitor laptops and smart devices.

Firewall rules between them follow least privilege from Chapter~\ref{ch:2.1}. Employees reach
servers on specific ports only: 443 for HTTPS, 445 for file sharing. Guest devices cannot reach any
internal VLAN at all. IoT devices can reach only the internet, never internal resources.

\begin{lstlisting}[language=Python, caption={Determining which traffic flows a VLAN scheme blocks}]
def blocked_flows(host_vlan, flows, allowed_pairs):
    """host_vlan: dict mapping host name -> VLAN id.
       flows: list of (src_host, dst_host) tuples.
       allowed_pairs: set of (src_vlan, dst_vlan) tuples permitted to
       communicate (same VLAN is always implicitly allowed).
       Returns sorted (src, dst) pairs that are blocked."""
    blocked = []
    for src, dst in flows:
        src_vlan, dst_vlan = host_vlan[src], host_vlan[dst]
        if src_vlan == dst_vlan:
            continue
        if (src_vlan, dst_vlan) not in allowed_pairs:
            blocked.append((src, dst))
    return sorted(blocked)
\end{lstlisting}

When a single switch port must carry traffic for multiple VLANs -- typically the uplink between a
switch and a router -- it runs in \textbf{trunk mode}, adding an 802.1Q tag to each frame
identifying its VLAN. The router strips the tag, applies its rules, and re-tags the outbound frame
for its destination VLAN. \textbf{Access ports}, where end-user devices plug in, carry exactly one
VLAN and never use tags at all -- the tagging is invisible to the endpoint.

\begin{misconception}
\textbf{``Setting up VLANs means attackers definitely cannot cross them.''} VLANs stop direct
layer-2 traffic, full stop. If the router or firewall rules governing inter-VLAN traffic are
misconfigured to permit everything, segmentation is defeated entirely. The firewall policy at the
routing boundary is what actually enforces the separation, not the VLAN tag itself.
\end{misconception}

\section{Traffic Isolation and Zero Trust}

VLANs solve the flat-network problem at a coarse grain, but a meaningful gap remains: within one
VLAN, every device can still freely reach every other device. If a single employee's laptop on the
``Employees'' VLAN is compromised, the attacker can scan and attack every other machine sharing
that VLAN without crossing a single boundary.

\begin{definition}{Zero Trust}
A security paradigm built on the principle ``never trust, always verify,'' a phrase coined by
analyst John Kindervag in 2010.\srccite{3.3}{1} Network access is never granted implicitly based on
a user's physical or logical location. Being plugged into an office Ethernet port, or connected to
corporate Wi-Fi, is treated as no more trustworthy than the public internet. Every request to
access a resource must be explicitly authenticated, authorized, and encrypted.
\end{definition}

Zero Trust rests on three principles, formalized under nearly these exact names in Microsoft's own
Zero Trust guidance.\srccite{3.3}{2} \textbf{Verify explicitly} means authenticating and authorizing
every request against every available signal -- user identity, device health, location, and anomaly
detection -- rather than trusting a session once it is established. \textbf{Use least-privilege
access} extends Chapter~\ref{ch:2.1}'s risk vocabulary with Just-In-Time and Just-Enough-Access
policies, granting a role only what it needs, only when it needs it. \textbf{Assume breach} means
designing as though an attacker is already inside. It calls for minimizing blast radius through
fine-grained segmentation, encrypting every session regardless of whether it is ``internal,'' and
monitoring continuously for anomalies rather than trusting that perimeter defenses already caught
everything.

\begin{definition}{Microsegmentation}
A security technique that applies fine-grained policy down to the individual workload or device
level, rather than grouping hundreds of devices into one VLAN and letting them communicate freely.
\end{definition}

Under microsegmentation, five hundred employee laptops are not treated as one trusted group.
Policies enforced by software-defined networking controllers, or by endpoint host firewalls, can
forbid laptops from communicating with each other entirely, permitting each to reach only specific
central servers and the internet. At the physical switch level, the same idea is implemented with
\textbf{Private VLANs (PVLANs)}, which divide one VLAN's broadcast domain into sub-domains using
three port types. \textbf{Isolated ports} cannot reach any other port in the same PVLAN except a
promiscuous one, even though they share the same subnet. \textbf{Community ports} can reach other
ports in their own community group and promiscuous ports, but nothing else. \textbf{Promiscuous
ports} can reach everything; they are typically reserved for the router or default gateway, so
every host can still reach the internet and DNS.

\begin{figure}[h]
\centering
\begin{tikzpicture}[
  port/.style={draw=kickstartermaroon, thick, rounded corners, minimum width=2.5cm, minimum height=0.8cm, align=center, font=\scriptsize},
]
\node[port, fill=kickstartermaroon!25] (prom) at (3.2,1.6) {Promiscuous port\\(router / gateway)};
\node[port, fill=kickstartercream] (iso1) at (0,0) {Isolated port A};
\node[port, fill=kickstartercream] (iso2) at (0,-1.3) {Isolated port B};
\node[port, fill=white] (com1) at (6.4,0) {Community port X};
\node[port, fill=white] (com2) at (6.4,-1.3) {Community port Y};
\draw[-{Latex[length=2mm]}, kickstartermaroon] (iso1) -- (prom);
\draw[-{Latex[length=2mm]}, kickstartermaroon] (iso2) -- (prom);
\draw[-{Latex[length=2mm]}, kickstartermaroon] (com1) -- (prom);
\draw[-{Latex[length=2mm]}, kickstartermaroon] (com2) -- (prom);
\draw[-{Latex[length=2mm]}, kickstartermaroon] (com1) -- (com2);
\draw[gray, dashed] (iso1) -- (iso2) node[midway, right, font=\tiny, gray] {never};
\draw[gray, dashed] (iso1) -- (com1) node[midway, above, font=\tiny, gray, sloped] {never};
\end{tikzpicture}
\caption{Private VLAN port types: isolated ports reach only the promiscuous port; community ports
reach their own community and the promiscuous port; the promiscuous port reaches everything.}
\label{fig:3.3-pvlan}
\end{figure}

\begin{lstlisting}[language=Python, caption={Evaluating a PVLAN traffic request by port type}]
def pvlan_allowed(src_type, dst_type, same_community=True):
    """src_type, dst_type: "isolated", "community", or "promiscuous".
       same_community: whether src and dst are in the same community
       group (only relevant when both are "community").
       Returns True if traffic is permitted."""
    if src_type == "promiscuous" or dst_type == "promiscuous":
        return True
    if src_type == "isolated" or dst_type == "isolated":
        return False
    return same_community
\end{lstlisting}

\begin{example}[The same attack under two architectures]
An attacker compromises a reception-desk PC. Under traditional VLANs, that PC shares VLAN~10 with
HR, finance, and printers. Ransomware on the compromised machine scans the subnet, finds an
exploitable service, and encrypts fifty other machines within minutes -- the entire VLAN was one
blast radius.

Under Zero Trust with microsegmentation, the same PC sits on an isolated PVLAN port. When the
malware attempts to scan its local subnet, the switch simply drops the packets, because an isolated
port cannot reach any peer. To reach a server at all, the request must pass through a Zero Trust
proxy that verifies the machine's security posture. Because that posture is now compromised
(unusual outbound scanning detected), its access token is revoked. The blast radius is exactly one
machine, not fifty.
\end{example}

\begin{lstlisting}[language=Python, caption={A minimal Zero Trust access-proxy evaluation}]
def evaluate_access_request(user_role, device_posture, resource_sensitivity):
    """user_role: str. device_posture: "healthy" or "compromised".
       resource_sensitivity: "low", "medium", or "high".
       Returns "allow" or "deny"."""
    if device_posture == "compromised":
        return "deny"
    if resource_sensitivity == "high" and user_role not in ("admin", "server-owner"):
        return "deny"
    return "allow"
\end{lstlisting}

\begin{misconception}
\textbf{``Once we set up VLANs, we have achieved Zero Trust.''} VLANs are a tool for traditional,
coarse segmentation based on IP routing. Zero Trust requires continuous verification of identity
and device health for every request -- a flat VLAN, even a well-designed one, still relies on
implicit trust for any two devices sharing it.
\end{misconception}

\section{Physical Air Gaps}

VLANs and Zero Trust microsegmentation still assume some network path exists, even if policy
restricts it. An \textbf{air gap} removes that path entirely: no network connection, wired or
wireless, links the protected system to any other network at all. It is segmentation's most
extreme case, the limit every other technique in this chapter approaches but does not reach.

\begin{definition}{Air gap}
A physical isolation strategy in which a system or network has no direct or indirect network
connection -- wired or wireless -- to any other network, including the internet. Data can cross an
air gap only through physical media or a human operator, never through a wire or a radio signal.
\end{definition}

Air gaps protect the systems with the least tolerance for compromise: nuclear-plant control
systems, classified military networks, and industrial control systems that run physical machinery.
An air gap defeats every network-based technique earlier in this chapter -- phishing, lateral
movement, a compromised Zero Trust proxy -- because no network path exists for any of them to use.

\begin{example}[Stuxnet: crossing an air gap without a network]
In 2010, researchers identified a worm, later named Stuxnet, running inside the centrifuge-control
network at Iran's Natanz uranium-enrichment facility -- a network with no direct connection to the
internet.\srccite{3.3}{3} Stuxnet did not cross that air gap over a wire. It spread through
infected USB drives, carried in by contractors whose own computers had been infected first, and
once inside it searched specifically for the industrial-control software running the centrifuges.
It then issued destructive commands to the centrifuges while feeding plant operators falsified
sensor readings showing normal operation, ultimately damaging close to a thousand of
them.\srccite{3.3}{3} An air gap eliminates the network path into a system. It does not eliminate
the person who still has to carry software onto it.
\end{example}

\begin{misconception}
\textbf{``An air-gapped system cannot be hacked.''} An air gap removes the network attack surface,
not every attack surface. Removable media, a compromised update process, or a person with physical
access remain valid paths in -- Stuxnet reached Natanz's isolated network without ever touching the
internet.
\end{misconception}

\section*{Review Questions}
\begin{reviewquestions}
\item An attacker compromises a guest laptop on an isolated guest VLAN with no inter-VLAN routing
rule to the server VLAN. Explain, step by step, why the attacker cannot reach the internal file
server even after fully compromising the guest laptop.
\item Explain the difference between an access port and a trunk port, and explain why the 802.1Q
tag used on a trunk port would break things if it were mistakenly applied to an access port.
\item Two employee laptops sit on the same VLAN with no microsegmentation. Explain why a
traditional VLAN boundary does nothing to stop one compromised laptop from directly attacking the
other, and describe the specific control that would.
\item A PVLAN has one isolated port and one community port in the same PVLAN. Explain whether
traffic between them is permitted, and explain what would need to change for them to communicate.
\item Explain why Zero Trust treats a device's physical connection to the corporate Wi-Fi network
as no more trustworthy than a connection from the public internet, and explain how this differs
from the assumption underlying traditional VLAN-based segmentation.
\item In the reception-desk PC example, ransomware on the compromised machine cannot even
enumerate other hosts on its own subnet under the Zero Trust / PVLAN architecture. Explain
mechanically why the scan itself fails, referencing the specific port type the reception PC sits
on.
\item A network architect proposes replacing VLANs entirely with microsegmentation, arguing VLANs
are now obsolete. Using the ``walls inside one building'' analogy from this chapter, explain why
VLANs and microsegmentation address different granularities of the same problem rather than one
simply superseding the other.
\item An air-gapped industrial control network has no wired or wireless connection to any other
network. Using the Stuxnet case, explain how an attacker still compromised it, and explain what
that reveals about the limits of segmentation as a defense.
\end{reviewquestions}

\begin{answerkey}
\item VLANs are logical broadcast domains created in switch software, so traffic broadcast on the
guest VLAN never physically reaches devices on the server VLAN unless a router or layer-3 switch
explicitly forwards it. With no inter-VLAN routing rule permitting that path, there is simply no
layer-2 route from the compromised guest laptop to the file server. Reaching it would require
traffic to leave through a router and be explicitly allowed by a firewall rule, neither of which
exists here. The compromise stays confined to the guest VLAN regardless of what the attacker does
on the laptop itself.
\answerlabel{3.3.A}
\item An access port carries exactly one VLAN's traffic and is what end-user devices plug into,
with no tagging visible to the endpoint at all. A trunk port carries multiple VLANs over a single
physical link by adding an 802.1Q tag identifying each frame's VLAN, typically used on the uplink
between a switch and a router. If the tag were mistakenly applied to an access port, the connected
device -- which expects untagged Ethernet frames -- would receive frames it cannot parse correctly.
It has no VLAN-aware network stack expecting an 802.1Q header.
\answerlabel{3.3.A}
\item A VLAN boundary only separates \emph{different} VLANs from each other. It does nothing to
restrict traffic between two devices that share the same VLAN, since a VLAN is one broadcast domain
in which every member is implicitly free to reach every other member. One compromised laptop can
therefore scan and attack the other directly. Microsegmentation -- enforced at the host level
through software-defined networking policy or, at the physical switch level, through Private VLANs
-- is the specific control that would stop this, by denying communication between individual hosts
even though they share a subnet.
\answerlabel{3.3.B}
\item Traffic between an isolated port and a community port in the same PVLAN is not permitted: an
isolated port can reach only a promiscuous port, never any other port in the PVLAN regardless of
its type. For them to communicate, at least one of the two would need to be reconfigured as a
promiscuous port (which can reach everything) or the isolated port would need to be changed to a
community port within the same community group as the other.
\answerlabel{3.3.A}
\item Zero Trust rests on the principle ``never trust, always verify'': a device's physical or
logical location is never treated as evidence of trustworthiness. Being plugged into corporate
Wi-Fi grants no more implicit access than a connection from the public internet would -- every
request must still be explicitly authenticated, authorized, and encrypted. Traditional VLAN-based
segmentation assumes the opposite: once a device is placed on a given VLAN, it is implicitly trusted
to communicate freely with every other device sharing that VLAN, with no per-request verification
at all.
\answerlabel{3.3.B}
\item The reception PC sits on an isolated PVLAN port, and an isolated port can reach only the
promiscuous port -- never any other port in the same PVLAN, including other machines on the exact
same subnet. When the ransomware attempts to scan its local subnet, every packet it sends toward
another host's IP is dropped at the switch itself, because the switch enforces the isolated port's
restriction in hardware regardless of what the infected machine's software attempts. The scan
produces no responses, because the probe packets never leave the switch in the first place.
\answerlabel{3.3.A}
\item VLANs solve the coarse-grained problem of separating large groups from each other -- HR from
Engineering from the guest lobby -- the way interior walls separate departments in a building.
Microsegmentation solves a finer-grained problem VLANs were never designed for: restricting
communication between individual devices \emph{within} one of those walled-off groups. It works the
way a locked office door restricts movement within a single department that already has its own
wall. Removing the walls (VLANs) and relying on microsegmentation alone would mean every device
authenticates individually, which Zero Trust supports. But coarse routing and broadcast-domain
separation still have to happen somewhere. The two techniques operate at different layers of the
same defense-in-depth strategy, not as competing replacements for each other.
\answerlabel{3.3.A}
\item Stuxnet reached Natanz's air-gapped centrifuge network without any network connection at
all, by spreading through infected USB drives carried in by people. An air gap closes only the
network path into a system; it does nothing to close a path that runs through removable media or a
person with physical access. Every segmentation technique in this chapter -- VLANs, Zero Trust, and
the air gap itself -- narrows or removes a \emph{network} path, but none of them substitutes for
controlling what physical media, and what people, are allowed to touch a protected system.
\answerlabel{3.3.A}
\end{answerkey}

\section*{Sources}
\begin{enumerate}[label=\lbrack\arabic*\rbrack, leftmargin=2.2em]
\item \srcitem{3.3}{1}John Kindervag, ``No More Chewy Centers: Introducing The Zero Trust Model
Of Information Security,'' Forrester Research, November 2010 -- the report that coined ``Zero
Trust'' and ``never trust, always verify.''
\item \srcitem{3.3}{2}Microsoft, ``Zero Trust security in Azure,'' Microsoft Learn -- states the
Verify explicitly / Use least privilege access / Assume breach principles as Microsoft's own
Zero Trust guidance.
\item \srcitem{3.3}{3}Kim Zetter, ``An Unprecedented Look at Stuxnet, the World's First Digital
Weapon,'' \emph{Wired}, November 3, 2014.
\end{enumerate}
